% Paper 9 (garch-x-vix) — R1 Revision: Dual-Target Robustness Subsection
% Source: K1066 (experiments/k1066/)
% Prepared: 2026-05-17
% Placement: Insert after \subsection{VIX vs.\ Macroeconomic Variables} in \section{Robustness}
% Integration notes:
%   - Add \label{sec:proxy_oc} cross-ref to Limitations (replace "future work" sentence)
%   - Limitations revision draft at bottom of this file
%   - No new figure required; Table alone is sufficient
%   - K1066 data: SPY OHLC + VIX, yfinance, 2005-01-04–2026-04-10, OOS 2019-01-01 (n=1828)

%=================================================================
\subsection{Alternative Volatility Proxy: Open-to-Close Returns}
\label{sec:proxy_oc}
%=================================================================

Our primary evaluation uses $r_t^2$ (close-to-close return squared) as the volatility proxy.
A potential concern is that models fitted on close-to-close returns ($r_{\mathrm{close}}$)
may align with the $r_{\mathrm{close}}^2$ proxy by construction, inflating the apparent advantage
of A4f when both the model and the evaluation criterion use the same return definition.

To test for proxy sensitivity, we estimate two additional models---GJR-GARCH and A4f
using open-to-close returns $r_{\mathrm{oc},t} = \log(\text{close}_t) - \log(\text{open}_t)$
as the return input---and evaluate all four models against $r_{\mathrm{oc}}^2$ as the
volatility target.\footnote{Open-to-close returns capture intraday variance only;
overnight returns are excluded. The $r_{\mathrm{oc}}^2$ proxy is noisier than
five-minute realized variance, which \citet{patton2011} shows preserves QLIKE rankings
regardless of proxy noise. A five-minute RV robustness check remains future work, but
the daily open-to-close test provides a meaningful first cross-check on proxy invariance.}
The rolling OOS design mirrors the main analysis: 2000-day estimation window,
63-day refit frequency, OOS 2019--2026 ($n = 1{,}828$).

Table~\ref{tab:dual_target_dm} reports the Diebold-Mariano statistics
across model--proxy combinations.

\begin{table}[H]
\centering
\caption{DM tests across return definitions and volatility proxies}
\label{tab:dual_target_dm}
\begin{threeparttable}
\small
\begin{tabular}{@{}llcccc@{}}
\toprule
Model A & Model B & Proxy & QLIKE (A) & DM $t$ & Harvey sig. \\
\midrule
\multicolumn{6}{l}{\textit{Baseline (main results, close-to-close target)}} \\[2pt]
A4f$_\mathrm{close}$ & GJR$_\mathrm{close}$ & $r^2_\mathrm{close}$ & $-8.360$ & $+4.03$ & Yes \\[4pt]
\multicolumn{6}{l}{\textit{Alternative proxy (open-to-close target)}} \\[2pt]
A4f$_\mathrm{oc}$ & GJR$_\mathrm{oc}$    & $r^2_\mathrm{oc}$ & $-8.816$ & $+4.04$ & Yes \\
A4f$_\mathrm{oc}$ & GJR$_\mathrm{close}$ & $r^2_\mathrm{oc}$ & $-8.816$ & $+7.05$ & Yes \\
\bottomrule
\end{tabular}
\begin{tablenotes}
\footnotesize
\item Positive DM $t$: Model A achieves lower QLIKE (better forecast).
\item Harvey~(2016) threshold: $|t| > 3.0$. SPY, OOS 2019--2026;
      baseline row ($t=4.03$) from Table~\ref{tab:pairwise_dm}; open-to-close rows
      from K1066 (same OOS window, $n=1{,}828$, seed 42).
\item Subscripts: \textit{close} = model estimated on $r_\mathrm{close}$;
      \textit{oc} = model estimated on $r_\mathrm{oc}$.
\end{tablenotes}
\end{threeparttable}
\end{table}

Three findings emerge from Table~\ref{tab:dual_target_dm}.
First, A4f$_\mathrm{oc}$ significantly outperforms GJR$_\mathrm{oc}$ on the open-to-close
target ($t = 4.04$, Harvey pass), confirming that the multiplicative VIX$^2$ structure
delivers statistically meaningful gains even when both the model and the benchmark
use open-to-close returns.
Second, the matched-model DM statistics are nearly identical across proxy definitions:
$t = 4.03$ for the close-to-close pair and $t = 4.04$ for the open-to-close pair.
This parallel strengthens the interpretation that the VIX$^2$ component adds genuine
forecast information rather than exploiting a mechanical alignment between the return
definition and the proxy.
Third, the cross-definition test---A4f$_\mathrm{oc}$ versus GJR$_\mathrm{close}$ evaluated
on $r^2_\mathrm{oc}$---yields $t = +7.05$, the largest DM statistic in our analysis.
This comparison isolates the combined gain from both switching to an open-to-close
model and incorporating VIX$^2$: relative to the plain close-to-close GJR baseline,
the open-to-close A4f achieves a seven-standard-deviation improvement in forecast
accuracy on the intraday variance target.

Sub-period stability further supports proxy invariance. A4f$_\mathrm{oc}$ achieves
lower QLIKE than GJR$_\mathrm{oc}$ in all five non-overlapping sub-periods spanning
2019--2026 (5/5; binomial $p = 0.031$), with QLIKE improvements ranging from
0.3\%~to~2.8\%. The result confirms directional consistency across market regimes
(Pre-COVID, COVID crash, Post-COVID, rate hike cycle, 2024--2026 expansion).

In summary, the advantage of the multiplicative VIX$^2$ framework documented in
Section~\ref{sec:results} is robust to the choice of volatility proxy:
the model delivers Harvey-significant gains whether the evaluation target is
close-to-close or intraday open-to-close variance, and the gains are quantitatively
similar in both settings.


% ---------------------------------------------------------------
% LIMITATIONS REVISION (replace existing "future work" sentence)
% ---------------------------------------------------------------
% Current text (line ~864 in main.tex):
%   "Since VIX and intraday realized variance have a non-trivial relationship,
%    the relative advantage of VIX-based models could change with proxy choice;
%    we leave this robustness check for future work with high-frequency data."
%
% Replacement:
%   "Since VIX and intraday realized variance have a non-trivial relationship,
%    the relative advantage of VIX-based models could change with proxy choice.
%    Section~\ref{sec:proxy_oc} provides a daily-frequency cross-check using
%    open-to-close returns as an alternative proxy, finding consistent Harvey-significant
%    gains ($t = 4.04$, 5/5 sub-periods); a five-minute realized variance check remains
%    for future work."
% ---------------------------------------------------------------
